\chapter{Markdown rendering}
\label{ch:markdown}

\begin{aside}
A README in the terminal, a help screen, a documentation
viewer --- all benefit from rendered Markdown. Plain text with
a few bold and italic spans goes a long way.
\end{aside}

\section{Parsing}

\maya{} apps can use \code{md4c}, \code{cmark}, or a small
hand-rolled parser. CommonMark is the standard.

\section{Block elements}

\begin{itemize}[leftmargin=*]
  \item Paragraph.
  \item Heading (1--6).
  \item Code block.
  \item Quote.
  \item Lists (ordered, unordered, nested).
  \item Tables.
  \item Horizontal rule.
\end{itemize}

Each maps to a \maya{} element with appropriate styling.

\section{Inline elements}

\begin{itemize}[leftmargin=*]
  \item Bold, italic, strikethrough.
  \item Inline code.
  \item Link (OSC-8 hyperlinks).
  \item Image (sometimes inline image protocol).
\end{itemize}

\section{Headings}

H1 might be bold + large + with a horizontal rule below.
H2: bold + colour. H3: bold only.

\section{Code blocks}

Render with monospace (already), optionally syntax-highlight
via tree-sitter or simple regex.

\section{Tables}

Use \maya{}'s table widget. Align by header marker (\code{:--},
\code{:-:}, \code{--:}).

\section{Reflow}

Wrap paragraphs to the terminal width. Reflow on resize.

\section{Images}

Inline images render via the protocols in
Chapter~\ref{ch:images}. For unsupported terminals, show
alt-text or a placeholder.

\section{Recap}

\begin{recap}
\begin{itemize}[leftmargin=*]
  \item Parse CommonMark with an off-the-shelf lib.
  \item Map blocks to \maya{} elements.
  \item Inline formatting via styles.
  \item Reflow on resize.
\end{itemize}
\end{recap}

\section*{Workshop}

\begin{enumerate}[leftmargin=*]
  \item Render a README in your app's help screen.
  \item Add syntax highlighting to code blocks.
  \item Support inline images via the half-block fallback.
\end{enumerate}
